%%--------------------------------------------------------
%% example.tex
%%
%%--------------------------------------------------------
\documentclass[12pt]{article}
% \input{prefix}%
\usepackage[sexy]{styles/evan}

\providecommand{\BibLatexMode}[1]{}
\providecommand{\BibTexMode}[1]{}

\ifx\UseBibLatex\undefined%
  \renewcommand{\BibLatexMode}[1]{}
  \renewcommand{\BibTexMode}[1]{#1}
\else
  \renewcommand{\BibLatexMode}[1]{#1}
  \renewcommand{\BibTexMode}[1]{}
\fi

%%%

\BibLatexMode{%
   \usepackage[bibencoding=utf8,style=alphabetic,backend=biber]{biblatex}%
   \usepackage{my_biblatex}%
}

%%%

% \newtheorem{theorem}{Theorem}[section]


%%

\providecommand{\emphind}[1]{}%
\renewcommand{\emphind}[1]{\emph{#1}\index{#1}}

\definecolor{blue25emph}{rgb}{0, 0, 11}

\providecommand{\emphic}[2]{}
\renewcommand{\emphic}[2]{\textcolor{blue25emph}{%
      \textbf{\emph{#1}}}\index{#2}}

\providecommand{\emphi}[1]{}%
\renewcommand{\emphi}[1]{\emphic{#1}{#1}}

\definecolor{almostblack}{rgb}{0, 0, 0.3}

\providecommand{\emphw}[1]{}%
\renewcommand{\emphw}[1]{{\textcolor{almostblack}{\emph{#1}}}}%

\providecommand{\emphOnly}[1]{}%
\renewcommand{\emphOnly}[1]{\emph{\textcolor{blue25}{\textbf{#1}}}}
%%

\newcommand{\HLink}[2]{\hyperref[#2]{#1~\ref*{#2}}}
\newcommand{\HLinkSuffix}[3]{\hyperref[#2]{#1\ref*{#2}{#3}}}

\newcommand{\figlab}[1]{\label{fig:#1}}
\newcommand{\figref}[1]{\HLink{Figure}{fig:#1}}

\newcommand{\thmlab}[1]{{\label{theo:#1}}}
\newcommand{\thmref}[1]{\HLink{Theorem}{theo:#1}}

\newcommand{\remlab}[1]{\label{rem:#1}}
\newcommand{\remref}[1]{\HLink{Remark}{rem:#1}}%

\newcommand{\corlab}[1]{\label{cor:#1}}
\newcommand{\corref}[1]{\HLink{Corollary}{cor:#1}}%

\providecommand{\deflab}[1]{}
\renewcommand{\deflab}[1]{\label{def:#1}}
\newcommand{\defref}[1]{\HLink{Definition}{def:#1}}

\newcommand{\lemlab}[1]{\label{lemma:#1}}
\newcommand{\lemref}[1]{\HLink{Lemma}{lemma:#1}}%

\providecommand{\eqlab}[1]{}%
\renewcommand{\eqlab}[1]{\label{equation:#1}}
\newcommand{\Eqref}[1]{\HLinkSuffix{Eq.~(}{equation:#1}{)}}


\newcommand{\MiraThanks}[1]{%
   \thanks{%
      School of Mathematical Sciences; %
      Queen Mary, University of London; %
      Mile End Road; %
      London, E14NS, UK; %
      \href{m.shamis@qmul.ac.uk}{m.shamis@qmul.ac.uk}; %
      \url{https://www.qmul.ac.uk/maths/profiles/shamiss.html}. %
   #1%
   }%
}


% Bib latex stuff
\BibLatexMode{%
   \usepackage[bibencoding=utf8,style=alphabetic,backend=biber]{biblatex}%
   \usepackage{my_biblatex}%
}

\BibLatexMode{%
   \bibliography{template}
}

\begin{document}

\title{Notes on Complex Variables MTH 5103 2024-2025}

\author{%
   Emomalijon Murodzoda%
   %
   \and%
   %
   Dr Mira Shamis%
   \MiraThanks{\emph{ASK HER IF YOU HAVE ANY QUESTIONS!!!}}
}

\date{\today}

\maketitle

\begin{abstract}
    %%% DELETE START
    One of the most annoying things about \LaTeX, and there are a lot,
    is how hard it is to just start writing math papers. This is an
    empty template intended to enable one to just start writing a
    paper. It uses my personal style, but might be useful for other people. It uses BibLatex by default. Search \texttt{DELETE} in the file to know which portion to delete when you start writing.
    %%% DELETE END
\end{abstract}


%%%%%%%%%%%%%%%%%%%%%%%%%%%%%%%%%%%%%%%%%%%%%%%%%%%%%%
%%%%%%%%%%%%%%%%%%%%%%%%%%%%%%%%%%%%%%%%%%%%%%%%%%%%%%

\section{Introduction}

%%% DELETE START 

%%%%%%%%%%%%%%%%%%%%%%%%%%%%%%%%%%%%%%%%%%%%%%%%
%%% TO USE TEMPLATE DELETE EVERYTHING FROM HERE...


All the \LaTeX{} definitions/macros/etc are in the \texttt{prefix.tex} file. This makes this file a bit cleaner.


\begin{theorem}
    \thmlab{first}%
    %
    Bla bla.
\end{theorem}
\begin{proof}
    This is the proof.
    \[
    \sqrt{2} + \sqrt{2} = 2 \sqrt{2}.
    \]
\end{proof}

% \james{It is extremely useful to leave comments in the text for your coauthors.}
% \ford{And the other author might reply...}
% \james{But one can still have the last word. }

See \thmref{first}.
\begin{lemma}
    \lemlab{second}.
    %
    This is a lemma.
\end{lemma}

See \lemref{second}.

\begin{definition}
    \deflab{number}
    A number is a \emphi{number}.
\end{definition}

See also \defref{number}.

\begin{remark}
    \remlab{useful}
    Remarks are useful sometime.
\end{remark}

Was \remref{useful} useful?

\section{Second section}


Do not forget to cite some irrelevant papers \cite{k-spda-10}.

%------------------------------------------------------------------
%------------------------------------------------------------------

\section{More stuff}
Some math:
% \begin{equation*}
%     \Var{X}%
%     =%
%     \Ex{(X - \Ex{X})^2}.
% \end{equation*}

Some definitions.
\begin{definition}
    \deflab{prime}%
    %
    An integer number $p > 1$ is \emphi{prime} if it is divisible
    only by $1$ or itself.
\end{definition}

A theorem:
\begin{theorem}
    \thmlab{main}% Label for the theorm
    %
    The number of primes is unbounded.
\end{theorem}
\begin{proof}
    Assume for the sake of contradiction that the number of primes is
    finite, say $k$, and let $1 < p_1< p_2 < \cdots < p_k$ be these
    primes. Observe that $N = p_1 \cdot p_2 \cdots p_k +1 > 1$ is
    indivisible by $p_1, \ldots, p_k$, and is larger than all these
    numbers. Thus, $N$ must be prime. A contradiction.
\end{proof}

\defref{prime} and \thmref{main} were both known to the Greeks. Well,
to some of the Greeks.

\bigskip

I prefer to use \texttt{enumitem} for creating lists:
% \begin{compactenumI}
%     \smallskip%
%     \item One can create more compact lists.

%     \smallskip%
%     \item There is more control over labels.
% \end{compactenumI}
% \smallskip%
% And another good reason is because:
% \begin{compactenumI}[resume]
%     \smallskip%
%     \item One can resume the numbering.
% \end{compactenumI}

\bigskip%

It is usually a good idea to let \LaTeX{} do its thing. Do not use
\textbackslash\textbackslash{} to end lines (common mistakes for
\LaTeX{} beginners). End a paragraph by having an empty line. If you
want a big space between two paragraphs, just put
\texttt{\textbackslash{big{}skip}} in a line on its own between th two
paragraphs.

\bigskip

Just like that.


\paragraph{Paragraphs.}

I like to title my paragraphs so that people know what the paragraph
is about. This is a personal style thing, as area lot of other writing
stuff. Follow what you like. More sectioning commands follow. If you
comment out the
\texttt{\textbackslash{}def\textbackslash{}UseBibLatex{1}} then the
system would use bibtex instead.

\subsection{A subsection}

\subsection{A lemma example}

\begin{lemma}
    This is a lemma.
\end{lemma}


%------------------------------------------------------------------
%------------------------------------------------------------------

\section{Bibliography}

Nowadays, I like to use \texttt{biblatex}, but it is somewhat more
painful to use than \texttt{bibtex}. The big advantage of
\texttt{biblatex} that it is highly configurable if you are willing to
spend the energy to learn it. It does much better work than
\texttt{bibtex}. I highly recommend getting your bibtex entry from
DBLP, since they have the \texttt{doi} information -- this implies
that you get a link to the paper in \texttt{biblatex} or
\texttt{bibtex} if you use the right style.

Anyway, here is an example of a citation \cite{k-spda-10}.





%%% DELETE END

\section*{Acknowledgments}

%%% DELETE START
Thank everyone.
%%% DELETE END

%-------------------------------------------------------------------------

\BibTexMode{%
   \bibliographystyle{alpha}
   \bibliography{template}
}%
\BibLatexMode{\printbibliography}

\end{document}


%--------------------------------------------------------
%
% x.tex - end of file
%--------------------------------------------------------

